
\section{Masse und Gewichtskraft}

	\begin{tabular}{cc}
	\parbox{3cm}{
		\begin{emphbox}
			\begin{align*}
				m &= \rho \cdot V \\
				F_G &= m \cdot g = \rho \cdot V \cdot g \\
				F_G &= m {\cdot} g = \rho {\cdot} V {\cdot} g
			\end{align*}
		\end{emphbox}
		} & 
	\parbox{3cm}{
		\begin{symbolbox}
			$m$ : Masse \\
			$\rho$ : Dichte \\
			$V$: \text{Volumen} \\
			$F_G$: Gewichtskraft \\
			$g$: \text{: Fallbeschleunigung}
		\end{symbolbox}
	}
	\end{tabular}


\section{Bewegung}
Gleichförmig:

	\begin{tabular}{cc}
	\parbox{2.5cm}{
		\begin{emphbox}
			\begin{align*}
				v = \frac{s}{t}, [v] = \frac{m}{s}
			\end{align*}
		\end{emphbox}
	} & 
	\parbox{3.5cm}{
		\begin{symbolbox}
			$v$ : Geschwindigkeit \\
			$s$ : Weg, Weglänge \\
			$t$ : Zeit, Dauer
		\end{symbolbox}
	}
	\end{tabular}


Gleichmässig beschleunigt/verzögert

	\begin{tabular}{cc}
	\parbox{2.5cm}{
		\begin{emphbox}
			\begin{align*}
				\text{Beschleunigung} a = \frac{\Delta v}{\Delta t}, [a] = \frac{m}{s^2} \\
				\text{Geschwindigkeit} v = \frac{\Delta s}{\Delta t} \\
				v = v_0 + a \cdot t \\
				\text{Weg} s = v_0 \cdot t + \frac{1}{2} a \cdot t^2 \\
				\text{Bremsweg} s_b = \frac{{v_0}^2}{2 \cdot a}
			\end{align*}
		\end{emphbox}
	} & 
	\parbox{3.5cm}{
		\begin{symbolbox}
			$a$ : Beschleunigung / Verzögerung $(a<0)$ \\
			$v_o$ : Anfangsgeschwindigkeit
		\end{symbolbox}
	}
	\end{tabular}

Freier Fall

\begin{tabular}{cc}
	\parbox{2.5cm}{
		\begin{emphbox}
			\begin{align*}
				v = g \cdot t, [g] = \frac{m}{s^2} \\
				\text{Fallhöhe:} h = \frac{1}{2} g \cdot t^2
			\end{align*}
		\end{emphbox}
	} & 
	\parbox{3.5cm}{
		\begin{symbolbox}
			$g$ : Fallbeschleunigung ($g = 9.81 m/s^2$)
		\end{symbolbox}
	}
\end{tabular}

Kreisförmige Bewegung

\begin{tabular}{cc}
	\parbox{2.5cm}{
		\begin{emphbox}
%			\begin{tabular}{lll}
			\begin{align*}
				\omega &= 2 \pi n, [\omega] = rad \cdot s^-1 ? s^-1 \\
				v &= 2 \pi n \cdot r = \pi n \cdot d = \omega \cdot r \\
				n &= \frac{1}{T}, [n] = s^{-1}
			\end{align*}
		%	\end{tabular}
		\end{emphbox}
	} & 
	\parbox{3.5cm}{
		\begin{symbolbox}
			\omega : \text{Winkel-/Drehgeschw.} \\
			n : \text{Drehfrequenz/-zahl}  \\ 
			T : \text{Umlaufdauer} \\
			r : \text{Bahnradius} \\
			d : \text{Bahndurchmesser} \\
			v : \text{Umfangsgeschw.}
		\end{symbolbox}
	}
\end{tabular}



\section{Akustik}

\begin{sectionbox}
	\subsection{Mach}

\begin{tabular}{cc}

	\parbox{2.5cm}{
		\begin{emphbox}
			\begin{align*}
				\sin \mu &= \frac{c}{v} = \frac{1}{M}\\
				\mu &= \arcsin \frac{1}{M}
			\end{align*}
		\end{emphbox}
	} & 
	\parbox{3cm}{
		\begin{symbolbox}
			c: Schallgeschwindigkeit \\
			v: Objektgeschwindigkeit \\
			M: Machzahl \\
			$\mu$: Machwinkel
		\end{symbolbox}
	}
\end{tabular}

	\resizebox{5cm}{!}{
		
		\def\dotMarkRightAngle[size=#1](#2,#3,#4){%
			\draw ($(#3)!#1!(#2)$) -- 
			($($(#3)!#1!(#2)$)!#1!90:(#2)$) --
			($(#3)!#1!(#4)$);
			\path (#3) --node[circle,fill,inner sep=.5pt]{} ($($(#3)!#1!(#2)$)!#1!90:(#2)$);
		}
		
		\begin{tikzpicture}
			
			\coordinate (OO) at (0,0);				
			
			% Circle centers
			\coordinate (P3) at (1,0);				
			\coordinate (P2) at (2,0);				
			\coordinate (P1) at (3,0);				
			
			% Tangential points
			\coordinate (S3) at ({1 *(1-sin(30)*sin(30)},{1*sin(30)*cos(30)});				
			\coordinate (S2) at ({2 *(1-sin(30)*sin(30)},{2*sin(30)*cos(30)});				
			\coordinate (S1) at ({3 *(1-sin(30)*sin(30)},{3*sin(30)*cos(30)});				
			
			% Endpoint of Mach front
			\coordinate (SX) at (4,{4*tan(30)});				
			
			\draw [-Stealth,dashed] (P3) -- node[above] {} (S3);
			\draw [-Stealth,dashed] (P2) -- node[above] {} (S2);
			\draw [-Stealth] (P1) -- node[right] {Schallgeschwindigkeit c} (S1);
			
			\coordinate (O) at (4,0);
			%		\coordinate (C) at (0,3);
			%		\coordinate (D) at ($(C)!(P3)!(B)$);
			
			\draw [-Stealth] (O) -- node[below] {Objektgeschwindigkeit v} (OO);
			\draw [dashed] (OO) -- node[left] {Machwelle} (SX);
			
			\pic [draw, angle radius=3mm, angle eccentricity=1.2mm] {angle = P3--OO--S3};
			\node[above = 0.1cm and 0.8cm of OO] {$\mu$};
			
			\dotMarkRightAngle[size=6pt](OO,S1,P1);
			
			\draw (P3) [dashed] circle ({1*sin(30)});
			\draw (P2) [dashed] circle ({2*sin(30)});
			\draw (P1) circle ({3*sin(30)});
			
			
		\end{tikzpicture}
		
	}		
	
	
	$M < 1.0$: Subsonic \\
	$M = 1.0$: Transonic \\
	$M > 1.0$: Supersonic \\
	$M > 5.0$: Hypersonic
	
\end{sectionbox}
